% Generated draft: requires mathematical and source review.
\documentclass[11pt]{article}
\usepackage[margin=25mm]{geometry}
\usepackage{fontspec}
\setmainfont{Latin Modern Roman}
\usepackage{amsmath,amssymb,mathtools,unicode-math}
\setmathfont{Latin Modern Math}
\usepackage{enumitem,needspace,xurl,hyperref}
\setlength{\parindent}{0pt}
\setlength{\parskip}{5pt}
\setlength{\emergencystretch}{3em}
\newcommand{\N}{\mathbb{N}}
\newcommand{\Z}{\mathbb{Z}}
\newcommand{\Q}{\mathbb{Q}}
\newcommand{\R}{\mathbb{R}}
\newcommand{\C}{\mathbb{C}}
\newcommand{\F}{\mathbb{F}}
\newcommand{\A}{\mathbb{A}}
\newcommand{\PP}{\mathbb{P}}
\newcommand{\E}{\mathbb{E}}
\newcommand{\eps}{\varepsilon}
\newcommand{\dd}{\,\mathrm{d}}
\begin{document}
\section*{UnsolvedMath ID: 30006554}
\textbf{Source problem code:} catalogue entry
\subsection{Definitions and mathematical statement}
Fix a graph \(H\), and write \(h=v(H)=|V(H)|\). The exponent in the supplied question presupposes that \(h\) is finite. The source does not specify the graph conventions, the meaning or normalization of \(L_\infty\)-distance, the convention for counting copies, or the proposed local-cover parameter. Consequently, a fully specified mathematical conjecture cannot be recovered from this record alone.

To display the intended quantifiers, let \(D_H(G)\) denote the normalized distance from \(G\) to the class of \(H\)-free graphs intended by the source, and let \(C_H(G,S)\) denote its intended number of copies of \(H\) extending a vertex set \(S\subseteq V(G)\). These are placeholders for missing definitions, not newly specified graph parameters. In particular, the record does not determine whether \(H\)-free means exclusion as an ordinary subgraph or as an induced subgraph.

Under the interpretation that the little-\(o\) bound is uniform over all vertex sets of the indicated size, the first question has the following conditional formulation: for every fixed \(H\), does there exist a nonnegative integer \(a=a(H)\), independent of \(n\) and of the host graphs, such that for every sequence \((G_n)\) with \(|V(G_n)|=n\),
\[
D_H(G_n)\longrightarrow 0
\quad\Longleftrightarrow\quad
\frac{\displaystyle\max_{\substack{S\subseteq V(G_n)\\ |S|=a}} C_H(G_n,S)}{n^{h-a}}
\longrightarrow 0
\qquad(n\longrightarrow\infty)?
\]
The maximum is taken for \(n\geq a\). Nonnegativity is necessary to interpret \(a\) as a set cardinality; the record states only that \(a(H)\) is an integer and gives no further admissible range. Uniformity is an explicit interpretation of the supplied wording and requires source confirmation.

The second question asks whether the proposed local-cover parameter gives a valid value of \(a(H)\) in this equivalence. Its definition is absent: the supplied original passage refers to a following definition without supplying it. Thus neither that parameter nor its asserted relationship to the displayed equivalence can be stated precisely here. The record also does not specify whether ``the correct value'' means an admissible value or an optimal value.
\subsection{Short English statement}
The question asks whether closeness to being \(H\)-free, measured locally in an unspecified \(L_\infty\) sense, is equivalent to uniformly small numbers of extensions of every vertex set of a fixed size \(a(H)\). It also asks whether a proposed local-cover parameter supplies that size. Missing definitions prevent a complete formulation; the displayed equivalence records the intended asymptotic structure conditionally.
\subsection{Sources}
\begin{itemize}
\item Local registry: permanent UnsolvedMath ID \(30006554\), title \emph{The Correct Parameter for Local Graph Removal}, source problem code \texttt{catalogue entry}, in the supplied Oberwolfach Reports Open Problems collection. The project identifies \texttt{data/unsolvedmath/problems.json} and \texttt{catalog/problems.json} as local registries; this draft uses the supplied complete record without inspecting those files.
\item Supplied report citation: \emph{Combinatorics} (2026), Oberwolfach Reports, DOI 10.4171/owr/2026/1, pp.\ 5--79. Source URL and supplied literature reference, labeled ``Oberwolfach Report 1/2026'': \url{https://doi.org/10.4171/owr/2026/1}. No individual proposer is identified; the supplied proposed year is 2026.
\item Literature assessment dated exactly \texttt{2026-08-22}: the supplied curation labels the problem partially solved and reports that the proposed parameter is verified for cliques, connected bipartite graphs, and further unspecified special cases, while the general characterization remains conjectural. This is supplied dated curation, not an independently verified literature assessment. The record expressly requests independent verification of time-sensitive status and cited claims.
\item The registry labels the statement \texttt{corrected\_verified} and reports that the official report was checked and the question isolated with repaired notation. These are supplied source-verification claims; this draft does not independently verify them. The report is described as documenting mathematical workshops and their open-problem discussions.
\end{itemize}
\end{document}
